\section{Результаты}
Сборка и запуск:

\begin{lstlisting}[language={}]
cmake -S . -B build
cmake --build build
./build/lab3 -t 4 -n 1000 -s 5
\end{lstlisting}

Вывод программы. Каждый поток печатает свой TID, порядок строк каждый раз разный, это показывает, что потоки работают параллельно и порядок их выполнения определяет планировщик:

\lstinputlisting[language={}]{inc/run.txt}

\subsection*{Количество потоков}
Скрипт \texttt{show\_threads.sh 4 3000} запускает программу в фоне с 4 потоками на 3000 точках и через секунду выводит информацию о потоках стандартными средствами Linux. У всех потоков один PID, а LWP (это TID) разные и совпадают с выводом программы. NLWP и строка \texttt{Threads} равны 5: четыре рабочих потока и главный. Главный поток спит в \texttt{pthread\_join} (состояние \texttt{S}, 0\% CPU), рабочие потоки выполняются (\texttt{R}) и загружают каждый своё ядро на 100\%.

\lstinputlisting[language={},columns=fixed]{inc/threads.txt}

То же самое можно посмотреть командами \texttt{ps -eLf}, \texttt{top -H -p PID} или в \texttt{htop}.

\subsection*{Исследование ускорения и эффективности}
Ускорение $S_p = T_1 / T_p$, эффективность $E_p = S_p / p$, где $T_p$ --- время работы на $p$ потоках. Замеры сделаны скриптом \texttt{bench.sh}: для каждой пары $(n, p)$ три запуска, взято лучшее время. Процессор Intel Core Ultra 9 285H: 16 ядер без Hyper-Threading, из них 6 производительных (P), 8 энергоэффективных (E) и 2 low-power E-ядра. Ubuntu запущена в WSL2.

\begin{center}
\begin{tabular}{|r|r|r|r|r|r|r|r|r|}
\hline
 & \multicolumn{2}{c|}{$n = 200$} & \multicolumn{2}{c|}{$n = 500$} & \multicolumn{2}{c|}{$n = 1000$} & \multicolumn{2}{c|}{$n = 2000$} \\
\hline
$p$ & $T$, мс & $S$ & $T$, мс & $S$ & $T$, мс & $S$ & $T$, мс & $S$ \\
\hline
1 & 1.659 & 1.00 & 22.40 & 1.00 & 183.5 & 1.00 & 1555.7 & 1.00 \\
2 & 0.810 & 2.05 & 11.36 & 1.97 & 95.9 & 1.91 & 800.1 & 1.94 \\
4 & 0.616 & 2.69 & 6.67 & 3.36 & 47.7 & 3.84 & 433.1 & 3.59 \\
8 & 0.573 & 2.90 & 3.81 & 5.88 & 27.1 & 6.78 & 246.8 & 6.30 \\
12 & 0.591 & 2.81 & 3.43 & 6.54 & 20.6 & 8.92 & 181.5 & 8.57 \\
16 & 0.701 & 2.37 & 2.89 & 7.76 & 18.0 & 10.22 & 171.2 & 9.09 \\
24 & 0.692 & 2.40 & 3.43 & 6.54 & 18.2 & 10.07 & 175.9 & 8.84 \\
32 & 0.849 & 1.95 & 3.11 & 7.21 & 18.1 & 10.16 & 176.7 & 8.80 \\
\hline
\end{tabular}
\end{center}

\begin{center}
\begin{tikzpicture}
\begin{axis}[
  width=0.9\textwidth, height=7.5cm,
  xlabel={Количество потоков $p$}, ylabel={Ускорение $S_p$},
  xmin=0, xmax=33, ymin=0, ymax=17,
  xtick={1,2,4,8,12,16,24,32},
  grid=major, legend pos=north east, legend cell align=left,
]
\addplot[black, dashed, domain=1:16] {x};
\addlegendentry{идеальное $S = p$}
\addplot[blue, mark=*] coordinates {(1,1.00) (2,2.05) (4,2.69) (8,2.90) (12,2.81) (16,2.37) (24,2.40) (32,1.95)};
\addlegendentry{$n = 200$}
\addplot[red, mark=square*] coordinates {(1,1.00) (2,1.97) (4,3.36) (8,5.88) (12,6.54) (16,7.76) (24,6.54) (32,7.21)};
\addlegendentry{$n = 500$}
\addplot[green!50!black, mark=triangle*] coordinates {(1,1.00) (2,1.91) (4,3.84) (8,6.78) (12,8.92) (16,10.22) (24,10.07) (32,10.16)};
\addlegendentry{$n = 1000$}
\addplot[violet, mark=diamond*] coordinates {(1,1.00) (2,1.94) (4,3.59) (8,6.30) (12,8.57) (16,9.09) (24,8.84) (32,8.80)};
\addlegendentry{$n = 2000$}
\end{axis}
\end{tikzpicture}

Рис. 1. Зависимость ускорения от количества потоков
\end{center}

\begin{center}
\begin{tikzpicture}
\begin{axis}[
  width=0.9\textwidth, height=7.5cm,
  xlabel={Количество потоков $p$}, ylabel={Эффективность $E_p$},
  xmin=0, xmax=33, ymin=0, ymax=1.1,
  xtick={1,2,4,8,12,16,24,32},
  grid=major, legend pos=north east, legend cell align=left,
]
\addplot[blue, mark=*] coordinates {(1,1.00) (2,1.02) (4,0.67) (8,0.36) (12,0.23) (16,0.15) (24,0.10) (32,0.06)};
\addlegendentry{$n = 200$}
\addplot[red, mark=square*] coordinates {(1,1.00) (2,0.99) (4,0.84) (8,0.73) (12,0.54) (16,0.49) (24,0.27) (32,0.23)};
\addlegendentry{$n = 500$}
\addplot[green!50!black, mark=triangle*] coordinates {(1,1.00) (2,0.96) (4,0.96) (8,0.85) (12,0.74) (16,0.64) (24,0.42) (32,0.32)};
\addlegendentry{$n = 1000$}
\addplot[violet, mark=diamond*] coordinates {(1,1.00) (2,0.97) (4,0.90) (8,0.79) (12,0.71) (16,0.57) (24,0.37) (32,0.28)};
\addlegendentry{$n = 2000$}
\end{axis}
\end{tikzpicture}

Рис. 2. Зависимость эффективности от количества потоков
\end{center}

{\bfseries Объяснение результатов.}

\begin{enumerate}
\item При маленьких данных ($n = 200$, около 1.6 мс работы) ускорение быстро упирается примерно в 3 раза, а после 8 потоков время начинает расти. Создание каждого потока (выделение стека через \texttt{mmap}, \texttt{clone3}) и его join стоят десятки микросекунд, и при большом числе потоков эти накладные расходы сравнимы с самой работой. По закону Амдала $S_p = \frac{1}{f + (1-f)/p}$ ускорение ограничено долей $f$ последовательной части, а здесь она большая.
\item С ростом $n$ ускорение растёт: объём работы растёт как $n^3$, а накладные расходы на потоки от $n$ почти не зависят, поэтому доля последовательной части уменьшается. На $n = 1000$ и $n = 2000$ эффективность до 4 потоков около 0.9--0.96, то есть почти идеальная. Потоки не обмениваются данными, массив точек только читается и помещается в кэш (2000 точек по 24 байта это 48 КБ), а мьютекс захватывается редко.
\item От 6 до 16 потоков ускорение продолжает расти, но эффективность падает до 0.6. Это связано с тем, что ядра у процессора разные: первые потоки попадают на быстрые P-ядра, следующие на более медленные E-ядра, последние на совсем слабые LP E-ядра. Кроме того, при нагрузке на все ядра ноутбук снижает частоту. В итоге максимум около 10 раз, а не 16.
\item Больше 16 потоков прироста нет: ядер всего 16, лишние потоки не работают параллельно, а делят ядра по времени, и добавляются переключения контекста. Ускорение остаётся на одном уровне, а эффективность падает примерно как $1/p$. Сильного замедления нет из-за динамической раздачи строк: поток, который реже получает процессор, просто берёт меньше строк.
\item На $n = 2000$ максимальное ускорение (9.1) немного меньше, чем на $n = 1000$ (10.2). Один запуск идёт больше секунды, процессор успевает нагреться и сбросить частоту. Также влияет то, что WSL2 работает в виртуальной машине, а Windows параллельно выполняет свои процессы. Значение $E > 1$ на $n = 200$, $p = 2$ это погрешность измерения очень маленького времени.
\end{enumerate}
